% ================================== Seitenränder ==================================

\KOMAoptions{DIV=calc}

% ================================== Seitenränder ==================================
% ================================== Metadaten ==================================

\newcommand{\workTitle}{Dokumentation}
\newcommand{\subTitle}{B-GOPB01-XX3-N01}
\newcommand{\specialization}{Grundlagen der objektorientierten Programmierung}
\newcommand{\studentFirstName}{Mücahid Emin}
\newcommand{\studentLastName}{Tomakin}
\newcommand{\studentId}{30042741}
\newcommand{\advisorPreTitle}{Vorg.-Titel}
\newcommand{\advisorFirstName}{Vorname}
\newcommand{\advisorLastName}{Nachname}
\newcommand{\advisorPosTitle}{Nachg. Titel}
\newcommand{\assessorPreTitle}{Vorg.-Titel}
\newcommand{\assessorFirstName}{Vorname}
\newcommand{\assessorLastName}{Nachname}
\newcommand{\assessorPosTitle}{Nachg. Titel}
\newcommand{\dateDay}{\the\day}
\newcommand{\dateMonth}{\the\month}
\newcommand{\dateYear}{\the\year}

% ================================== Metadaten ==================================
% ================================== Packages ==================================

% =========== Layout & Formatierung =====
\KOMAoptions{parskip=full} % Abstand zwischen Absätzen
\usepackage[left=3.5cm, right=3cm, bottom=3cm, top=3cm]{geometry} % Seitenränder
\usepackage[onehalfspacing]{setspace} % Zeilenabstand
\usepackage[T1]{fontenc} % Korrekte Zeichencodierung
\usepackage[utf8]{inputenc} % Eingabecodierung (nicht mehr notwendig bei LuaLaTeX/XeLaTeX)
\usepackage[ngerman]{babel} % Deutsche Sprache für Trennungen, Datum, Kapitel usw.
\usepackage{microtype} % Bessere Typografie (Silbentrennung, Randausgleich etc.)
\usepackage{enumitem} % Bessere Kontrolle über Listen
\usepackage{multicol} % Mehrspaltiger Text
\usepackage{blindtext} % Platzhaltertext
\usepackage{comment} % mehrere Zeilen auf einmal zu kommentieren
% =========== Layout & Formatierung =====
% =========== Mathematik & Symbole =====
\usepackage{amsmath} % Erweiterte Matheumgebung
\usepackage{amssymb} % Mathematische Symbole
\usepackage{mathtools} % Erweiterungen zu amsmath
\usepackage{fixmath} % Feste Mathe-Symbole
\usepackage{textcomp} % Zusätzliche Zeichen (z.B. °C)
\usepackage{fmtcount} % Römische Großbuchstaben
% =========== Mathematik & Symbole =====
% =========== Tabellen & Listen =====
\usepackage{listings}
\usepackage{longtable} % Tabellen über mehrere Seiten
\usepackage{multirow} % Zellen über mehrere Zeilen
\usepackage{booktabs} % Schönere Tabellen (z.B. \toprule)
\usepackage{tabularx} % Tabellen mit automatischer Spaltenbreite
\usepackage{enumerate} % Nummerierte Listen mit mehr Kontrolle
% =========== Tabellen & Listen =====
% =========== Grafiken & Diagramme =====
\usepackage{graphicx} % Bilder einfügen
\usepackage[export]{adjustbox} % Zusätzliche Optionen für includegraphics
\usepackage{float} % Exakte Platzierung von Abbildungen ([H])
\usepackage{scrhack} % Für KOMA-Script Kompatibilität mit float
\usepackage{subcaption} % Unterabbildungen (aktuelle Version von 'subfigure')
\usepackage{grffile} % Unterstützt komplexe Dateinamen
\graphicspath{{Figures/}} % Standardpfad für Abbildungen
\usepackage{animate} % Animationen in PDF
\usepackage[dvipsnames]{xcolor} % Erweiterte Farbauswahl
\usepackage{pgfplots} % Diagramme mit TikZ
\pgfplotsset{compat=1.10} % Kompatibilität für pgfplots
\usepgfplotslibrary{statistics} % Statistik in pgfplots
\usepackage{pgf-pie} % Kreisdiagramme
\usepackage{tikz} % Zeichnungen
\usetikzlibrary{shapes,arrows} % Zusätzliche TikZ Bibliotheken
\usepackage{smartdiagram} % Diagramme
\usepackage[european,nooldvoltagedirection]{circuitikz} % Schaltpläne
% =========== Grafiken & Diagramme =====
% =========== Quellcode & Listings =====
\usepackage{lstautogobble} % Quellcode einbinden
\usepackage[breakable]{tcolorbox} % Codeboxen
\usepackage{fancyvrb} % Erweiterte verbatim-Umgebung
%\usepackage[style=numeric,backend=biber]{biblatex} % numeric als Beispiel
% =========== Quellcode & Listings =====
% =========== Kopf-/Fußzeile & Seitenzahlen =====
\usepackage[automark]{scrlayer-scrpage} % KOMA-Script Kopf-/Fußzeilen
\usepackage{lastpage} % Zugriff auf letzte Seite
% =========== Kopf-/Fußzeile & Seitenzahlen =====
% =========== Hyperlinks & URLs =====
\usepackage{url} % URLs setzen
\usepackage{hyperref} % Hyperlinks
\hypersetup{
    unicode=false,
    pdftoolbar=true,
    pdfmenubar=true,
    pdffitwindow=false,
    pdfstartview={FitV},
    pdfpagelayout={SinglePage},
    pdftitle={\workTitle},
    pdfsubject={\subTitle},
    pdfauthor={\studentFirstName~\studentLastName},
    pdfkeywords={\workTitle,~\subTitle,~\specialization},
    pdfcreator={pdflatex},
    pdfproducer={LaTeX with hyperref},
    pdfnewwindow=true,
    colorlinks=false,   % Farbige Links aktivieren
    linkcolor=blue,    % Farbe für interne Links (z.B. Inhaltsverzeichnis)
    citecolor=blue,    % Farbe für Zitate
    filecolor=blue,    % Farbe für Dateilinks
    menucolor=black,   % Farbe für Menulinks
    urlcolor=blue,     % Farbe für URLs
    breaklinks=true,
    pdfborder={0 0 0} % Keine Ränder um Links
}
% =========== Hyperlinks & URLs =====
% =========== Abkürzungsverzeichnis =====
\usepackage[nohyperlinks]{acronym} % Abkürzungen
% =========== Abkürzungsverzeichnis =====
% =========== Weitere nützliche Pakete =====
%\usepackage{csquotes} % Intelligente Anführungszeichen (empfohlen mit biblatex)
% =========== Weitere nützliche Pakete =====

% ================================== Packages ==================================
% ================================== Verzeichnisnamen ==================================

\renewcommand{\contentsname}{Inhaltsverzeichnis}
\renewcommand{\bibname}{Literaturverzeichnis}
\renewcommand{\lstlistoflistingsname}{Codeverzeichnis}
\renewcommand{\listtablename}{Tabellenverzeichnis}
\renewcommand{\listfigurename}{Abbildungsverzeichnis}

% ================================== Verzeichnisnamen ==================================
% ================================== Literaturverzeichnis ==================================

%\addbibresource{literature/bibliography.bib}

% ================================== Literaturverzeichnis ==================================
% ================================== BlockDiagram Config ==================================

\tikzset{%
  block/.style    = {draw, thick, rectangle, minimum height = 3em, minimum width = 3em},
  sum/.style      = {draw, circle, node distance = 2cm}, % Adder
  input/.style    = {coordinate}, % Input
  output/.style   = {coordinate}, % Output
  mult/.style     = {draw, isosceles triangle, minimum height=1cm, minimum width =1cm}
}

% ================================== BlockDiagram Config ==================================
%  ================================== Listings  ==================================

% =========== Einstellungen für Listings ===========

\lstset{
    title=\lstname, % Titel des Listings ist der Dateiname
    frame=single, % Einzelne Umrandung des Codes
    framerule=0pt, % Keine Dicke der Umrandung
    rulecolor=\color{lightgray}, % Graue Rahmenfarbe
    showspaces=false, % Keine Anzeige von Leerzeichen
    showstringspaces=false, % Keine Anzeige von Leerzeichen in Strings
    showtabs=false, % Keine Anzeige von Tabulatoren
    numbers=left, % Zeilennummern links
    numbersep=5pt, % Abstand der Zeilennummern
    numberstyle=\sffamily\tiny\color{gray}, % Stil der Zeilennummern
    stepnumber=1, % Schrittweite der Nummerierung
    breaklines=true, % Zeilenumbruch aktivieren
    autogobble=true, % Entfernt leere Präfixe
    linewidth=\textwidth, % Maximale Breite für Listings
    basicstyle=\sffamily\scriptsize, % Basisstil für den Code
    sensitive=true, % Case-sensitive Modus
    morekeywords={MyClass,MyFunction,def,class}, % Weitere Schlüsselwörter
    emph={MyClass,MyFunction,def,class}, % Besonders hervorzuhebende Schlüsselwörter
    emphstyle=\color{orange}\bfseries, % Stil für die hervorgehobenen Wörter
    commentstyle=\color{green}\itshape, % Grüne, kursive Kommentare
    keywordstyle=\color{blue}, % Blaue Schlüsselwörter
    stringstyle=\color{red}, % Rote Strings
    identifierstyle=\color{purple}, % Lila Bezeichner
    numberstyle=\tiny\color{gray} % Graue Zeilennummern
}
% =========== Definition des Input-Listings-Stils ===========
\lstdefinestyle{inputlistingstyle}{
    language=Python,
    frame=single,
    numbers=left,
    breaklines=true,
    basicstyle=\sffamily\scriptsize,
    captionpos=b, % Setzt die Caption oben
    abovecaptionskip=11pt, % Abstand über dem Titel
    belowcaptionskip=5pt % Abstand unter dem Titel
}
% =========== Definition der Codeboxen ===========
\newtcolorbox{codeblock}{
    colback=gray!5!white, % Hintergrundfarbe
    %colframe=gray!88!black, % Rahmenfarbe
    colframe=white, % Rahmenfarbe
    before skip=20pt, % Abstand vor der Box
    after skip=20pt, % Abstand nach der Box
    breakable, % Ermöglicht den Seitenumbruch innerhalb der Box
    top=-\tcboxrule, % Entfernt den oberen Rand
    bottom=-\tcboxrule % Entfernt den unteren Rand
}

%  ================================== Listings  ==================================
%  ================================== Color definition ================================== 

\definecolor{zB}{RGB}{0, 84, 150}
\definecolor{DB}{RGB}{0, 84, 150}
%\definecolor{DB}{RGB}{0, 0, 0}

%  ================================== Color definition ================================== 
% ================================== headline & titel space ==================================

\renewcommand*{\chapterheadstartvskip}{\vspace*{-.4cm}}
\renewcommand*{\chapterheadendvskip}{\vspace{.5cm}}
\setlength{\parindent}{0pt} % no indent at paragraph start

% ================================== headline & titel space ==================================
% ================================== Header and Footer ==================================

% Kopfzeile
\clearpairofpagestyles
\ohead{\subTitle} % Rechts in der Kopfzeile
\chead{} % Zentriert in der Kopfzeile
\ihead{} % Links in der Kopfzeile

% Fußzeile
\ofoot{\pagemark} % Rechts in der Fußzeile
\cfoot{} % Zentriert in der Fußzeile
\ifoot{\studentLastName~|~\studentFirstName} % Links in der Fußzeile

% Linien einstellen
\setkomafont{pageheadfoot}{\normalfont}
\setheadtopline{0.4pt}
\setfootbotline{0.4pt}
\pagestyle{scrheadings}

% ================================== Header and Footer ==================================
% ================================== Size of headings ==================================

\setkomafont{chapter}{\LARGE}
\setkomafont{section}{\Large}
\setkomafont{subsection}{\large}
\setkomafont{subsubsection}{\normalsize}
\setkomafont{paragraph}{\normalsize}
\setkomafont{subparagraph}{\small}

% ================================== Size of headings ==================================
% ================================== hyperref at the end ==================================

\hypersetup{
    unicode=false,
    pdftoolbar=true,
    pdfmenubar=true,
    pdffitwindow=false,
    pdfstartview={FitV},
    pdfpagelayout={SinglePage},
    pdftitle={\workTitle},
    pdfsubject={\subTitle},
    pdfauthor={\studentFirstName~\studentLastName},
    pdfkeywords={\workTitle,~\subTitle,~\specialization},
    pdfcreator={pdflatex},
    pdfproducer={LaTeX with hyperref},
    pdfnewwindow=true,
    colorlinks=false,
    linkcolor=black,
    citecolor=black,
    filecolor=black,
    menucolor=black,
    urlcolor=black,
    breaklinks=true
}

% ================================== hyperref at the end ==================================